\documentclass[../readings.tex]{subfiles}

\begin{document}

\subsection{Computational Complexity}
\label{sub:complexity}

Algorithmic cost measured in arithmetic operations (flops). Determines scalability and feasibility.

\subsubsection{Big-O Notation}

\addDef{Big-O Notation}{%
$f(n) = O(g(n))$ if $\exists C, n_0 > 0$ s.t. $|f(n)| \leq C|g(n)|$ for $n \geq n_0$.
\begin{itemize}
  \item $\Omega(g)$: Grows at least as fast.
  \item $\Theta(g)$: Grows at the same rate.
  \item $o(g)$: Grows strictly slower.
\end{itemize}
}
\label{def:big-o}

\addNote{%
\textbf{Typical Rates.}\enspace
$O(1) \ll O(\log n) \ll O(n) \ll O(n \log n) \ll O(n^2) \ll O(n^3) \ll O(2^n)$.
For $n=1000$: $O(n^2)$ is $\sim 10^6$ ops (ms); $O(n^3)$ is $\sim 10^9$ ops (s).
}

\addNote{%
\textbf{Properties.}\enspace
1. \textbf{Transitivity:} $f=O(g), g=O(h) \Rightarrow f=O(h)$.
2. \textbf{Sum Rule:} $O(f) + O(g) = O(\max(f, g))$.
3. \textbf{Product Rule:} $O(f) \cdot O(g) = O(fg)$.
4. \textbf{Constant Factors:} $O(cf) = O(f)$.
}

\addExcr{%
\begin{enumerate}
  \item Show $3n^2 + 100n + 5 = O(n^2)$.
  \item Compare growth: $n \log n$ vs $n^2$.
  \item Rank by cost: $5n^3+n, 100n^2, n^2 \log n, 2^{10}n^2$.
  \item Use Master Theorem: $T(n) = 2T(n/2) + n$.
\end{enumerate}
}

\subsubsection{Flop Counting}

Hardware-independent measure of work.

\addDef{Flop}{%
A single floating-point operation ($+, -, \times, \div$).
}
\label{def:flop}

\addThm{Basic Costs}{%
For $\bx, \by \in \fR^n$ and $\bA \in \fR^{m \times n}, \bB \in \fR^{n \times p}$:
\begin{itemize}
  \item \textbf{Dot product:} $2n$ flops.
  \item \textbf{Matrix-vector product:} $2mn$ flops.
  \item \textbf{Matrix-matrix product:} $2mnp$ flops ($2n^3$ for square).
  \item \textbf{Triangular solve:} $n^2$ flops.
\end{itemize}
}
\label{thm:flop-costs}

\addNote{%
\textbf{Note: The Memory Wall.}\enspace
In modern hardware, arithmetic is cheap but moving data (RAM to CPU) is expensive. An $O(n^3)$ algorithm with high cache locality (like blocked GEMM) can outperform an $O(n^2)$ algorithm with poor locality.
}

\addNote{%
\textbf{Benchmarks.}\enspace
\begin{itemize}
    \item \textbf{LU Factorization:} $\frac{2}{3}n^3$.
    \item \textbf{Cholesky:} $\frac{1}{3}n^3$.
    \item \textbf{QR (Householder):} $\frac{4}{3}n^3$.
\end{itemize}
}

\addExcr{%
\begin{enumerate}
  \item Verify $2mn$ flops for $\bA\bx$ by counting multiply-adds.
  \item Estimate wall-clock time for $10^9$ flops on a $10^9$ flop/s CPU.
  \item Python: Plot \texttt{np.dot(A, B)} time vs $n$ (log-log). Verify exponent $\approx 3$.
  \item Why might the empirical exponent be less than 3? (Cache hierarchy, vectorization).
\end{enumerate}
}

\end{document}
